\section{An exact cubic comparison for an entire local Hecke family}
\label{sec:cubic-family-hecke}

The two-point calculation AC-G38 can be strengthened without discarding
the original congruences between eigenpackets.  One additional
ramification hypothesis at coefficient three gives the complete
paired-family algebra, including those congruences.

\begin{theorem}[AC-G39: a whole-family cubic fiber product]
\label{AC-G39}
Let $\mathcal O$ be the ring of integers of a finite extension of
$\mathbb Q_\ell$ containing a primitive cube root $\omega$, and let
$\delta=\omega-1$.  Let $X$ be a nonempty finite set, let $\Gamma$
be profinite, and let $\chi:\Gamma\to\mu_3$ be continuous and
surjective.  Put $H=\ker\chi$.  For each $x\in X$, let
$V_x:\Gamma\to\operatorname{GL}_n(\mathcal O)$ be continuous.
Let $\mathcal F$ meet every conjugacy class of every finite quotient
of $\Gamma$.  Write $c_j(x,g)$ for the coefficient of $T^{n-j}$ in
$\det(T-V_x(g))$, and let
\[
 T_0=\mathcal O[\,(c_j(x,g))_{x\in X}:g\in\mathcal F,\ 1\le j\le n\,]
                    \subset\mathcal O^X.                    \tag{G39.1}
\]
Assume that $T_0$ is local.  Suppose that for one $x_0\in X$,
over an algebraic closure of the residue field,
\[
 \overline V_{x_0}^{\rm ss}=S\oplus\bigoplus_a\eta_a,
 \qquad \dim S>1,\quad \dim\eta_a=1,
 \qquad S|_H\text{ irreducible},                            \tag{G39.2}
\]
with $S$ occurring once.  Let $T'$ be the $\mathcal O$-algebra
in $\mathcal O^X\times\mathcal O^X$ generated by diagonal scalar
constants and the pairs
\[
           \bigl(c_j(\cdot,g),\chi(g)^j c_j(\cdot,g)\bigr),
                  \qquad g\in\mathcal F,\quad 1\le j\le n.
                                                               \tag{G39.3}
\]
If $\ell\ne3$, then $T'=T_0\times T_0$.
If $\ell=3$, assume in addition that there exists $\tau\in\Gamma$
such that
\[
              \chi(\tau)=\omega,
              \qquad V_x(\tau)=1\quad\text{for every }x\in X.
                                                               \tag{G39.4}
\]
Then
\[
 T'=T_0\times_{T_0/\delta T_0}T_0
   =\{(a,b)\in T_0\times T_0:b-a\in\delta T_0\}.             \tag{G39.5}
\]
In particular, $T'$ contains the diagonal copy of $T_0$ and
$\delta P_0,\delta P_1$, where $P_i$ are its rational branch
projectors.  If $r_0=\operatorname{rank}_{\mathcal O}T_0$, then
\[
 \operatorname{lenUnit}_3((T_0\oplus T_0)/T')=r_0/2.           \tag{G39.6}
\]
No bound on the original congruence ideals of $T_0$ is assumed.
\end{theorem}
\begin{proof}
\proofstep{1. Finite integral closure and common residual traces.}
Both algebras are $\mathcal O$-submodules of finite free
$\mathcal O$-modules.  Since $\mathcal O$ is a noetherian DVR,
they are finite free and closed for the valuation topology.
For the latter assertion, a submodule of $\mathcal O^d$ has,
after a change of basis, the form
$\pi^{a_1}\mathcal O\oplus\cdots\oplus
\pi^{a_r}\mathcal O\oplus0$; every displayed summand is closed.
Thus neither the definition by algebraic generators nor passage
to all group elements requires an additional completion.

For each $g\in\Gamma$, continuity and the density of
$\mathcal F$ put every vector $c_j(\cdot,g)$ in $T_0$ and
its corresponding pair in $T'$.  Indeed, modulo $\pi^m$ all
finitely many $V_x$ and $\chi$ factor through a common finite
quotient.  A Frobenius class in that quotient gives the desired
simultaneous congruence; closedness permits $m\to\infty$.

Let $k=\mathcal O/\pi\mathcal O$.  For each $x$, evaluation
followed by reduction is a surjection $\epsilon_x:T_0\to k$
because $T_0$ contains scalar constants.  Its kernel is maximal.
Locality of $T_0$ makes these kernels equal.  For $a\in T_0$,
$a-a(x)1$ belongs to $\ker\epsilon_x$, hence to every
$\ker\epsilon_y$.  Therefore
\[
                      a(x)\equiv a(y)\pmod\pi
                           \quad(x,y\in X).                  \tag{G39.7}
\]
This applies, in particular, to the traces at every group element.

\proofstep{2. The coordinate images and the universal upper inclusion.}
Both coordinate projections of $T'$ have image $T_0$.
For the second coordinate, multiplying a generator by
$\chi(g)^{-j}\in\mathcal O^\times$ recovers the corresponding
generator of $T_0$.  Moreover, each generator in (G39.3) has
coordinate difference in $\delta T_0$.  The set of pairs with
this property is an algebra: for two pairs $(a,b),(a',b')$,
\[
                  bb'-aa'=b(b'-a')+a'(b-a).
\]
Consequently
$T'\subset T_0\times_{T_0/\delta T_0}T_0$.

\proofstep{3. Separation away from three.}
Choose $g$ with $\chi(g)=\omega$.  The coset-trace argument
proved in AC-G38, applied to (G39.2), gives $h\in H$ with
$\operatorname{tr}V_{x_0}(gh)$ a unit.  By (G39.7), all
$a_x=\operatorname{tr}V_x(gh)$ have the same nonzero residue
$\bar a$.  The element
\[
                         t=((a_x)_x,(\omega a_x)_x)
\]
belongs to $T'$ by Step 1.  Suppose $\ell\ne3$.  The two
monic polynomials
\[
             p_0(Z)=\prod_{x\in X}(Z-a_x),\qquad
             p_1(Z)=\prod_{x\in X}(Z-\omega a_x)
\]
are relatively prime over $\mathcal O$: their resultant
is, up to sign, $\prod_{x,y}(a_x-\omega a_y)$, a unit
because $\bar a\ne0$ and $\bar\omega\ne1$.
The Sylvester matrix is therefore invertible over
$\mathcal O$, giving polynomials $u,v\in\mathcal O[Z]$
with $up_0+vp_1=1$.  Evaluating $vp_1$ at $t$ gives
$P_0=(1,0)\in T'$, and $P_1=1-P_0\in T'$.
The surjectivity of both coordinate projections now gives
$T'=T_0\times T_0$.

\proofstep{4. A constant branch separator at three.}
Suppose $\ell=3$ and (G39.4) holds.  Apply the same
coset-trace argument with $g=1$ to choose $h\in H$ for
which $a_{x_0}=\operatorname{tr}V_{x_0}(h)$ is a unit.
Again every $a_x=\operatorname{tr}V_x(h)$ is a unit.
The trace vectors at $h$ and $\tau h$, respectively, are
\[
                a=((a_x)_x,(a_x)_x),\qquad
                b=((a_x)_x,(\omega a_x)_x)\quad\text{in }T'.
                                                               \tag{G39.8}
\]
Here the equality $V_x(\tau h)=V_x(h)$ uses precisely
the common hypothesis (G39.4).

The element $a$ is invertible inside $T'$, not just in the
ambient product.  To see this, write
$p(Z)=\prod_x(Z-a_x)=c_0+Zq(Z)$.  Its constant coefficient
$c_0$ is a unit, and $p(a)=0$, so
$a^{-1}=-c_0^{-1}q(a)\in T'$.
Thus
\[
                 U=ba^{-1}=(1,\omega)\in T',\qquad
                 U-1=\delta P_1,\qquad
                 \omega-U=\delta P_0.                       \tag{G39.9}
\]
Surjectivity of the two coordinate maps now gives
$\delta T_0\times0\subset T'$ and $0\times\delta T_0\subset T'$.
For any $f\in T_0$, choose $(f,g)\in T'$.  Step 2 gives
$g-f=\delta z$ for some $z\in T_0$.  Subtracting
$(0,\delta z)\in T'$ gives $(f,f)\in T'$.
We have proved the reverse inclusion in (G39.5).

Finally $(a,b)\mapsto b-a\bmod\delta T_0$ identifies the
quotient in (G39.6) with $T_0/\delta T_0$.  The module
$T_0$ is free of rank $r_0$, and $\delta^2=-3\omega$.
Dividing its length by $v_{\mathcal O}(3)$ gives $r_0/2$.
\end{proof}

\begin{corollary}[AC-G39-lattice: the added cost beyond old congruences]
\label{AC-G39-lattice}
Let a full $\mathcal O$-lattice $M\subset W_0\oplus W_1$
be stable under the paired-family algebra of AC-G39, with
the two branches acting on their corresponding eigenspace
sums.  Put $N_i=P_iM$.  For a decomposition
$W_i=Y\oplus Y'$ let $Q:W_i\to Y$ be its projection.
Then
\[
\begin{split}
 \operatorname{lenUnit}_\ell\bigl(QN_i/(M\cap Y)\bigr)
 &=\operatorname{lenUnit}_\ell\bigl(QN_i/(N_i\cap Y)\bigr)
   +\operatorname{lenUnit}_\ell\bigl((N_i\cap Y)/(M\cap Y)\bigr),\\
 0\le\operatorname{lenUnit}_\ell\bigl((N_i\cap Y)/(M\cap Y)\bigr)
 &\le
 \begin{cases}0,&\ell\ne3,\\ \dim Y/2,&\ell=3.
 \end{cases}
\end{split}                                                   \tag{G39.10}
\]
The first term is the original projection/intersection
index inside the projected branch lattice $N_i$; it is
retained without an upper bound.
\end{corollary}
\begin{proof}
Since $P_i$ preserves $M$ away from three, there
$N_i=M\cap W_i$, giving equality of the two intersections.
At three, $\delta P_i$ preserves $M$, so
$\delta N_i\subset M\cap W_i\subset N_i$ and hence
\[
             \delta(N_i\cap Y)\subset M\cap Y\subset N_i\cap Y.
\]
The lattice $N_i\cap Y$ is free of rank $\dim Y$.
Its quotient by $M\cap Y$ has normalized length at most
$\dim Y/2$.  The inclusions
$M\cap Y\subset N_i\cap Y\subset QN_i$ give the exact
sequence whose length additivity is (G39.10).
\end{proof}

\begin{proposition}[AC-G39-necessity: locality alone cannot separate a family]
\label{AC-G39-necessity}
At coefficient three, (G39.4) cannot simply be omitted
from the family statement, even if $T_0$ is local and
separates the original eigensystems and every residual
constituent satisfies (G39.2).
\end{proposition}
\begin{proof}
Take $\mathcal O$ containing $\omega$ and a primitive
fourth root $i$, and let $\Gamma=C_3\times Q_8$ with
$\chi$ the first-factor character.  Let $V$ be trivial
on $C_3$ and the following two-dimensional representation
on the quaternion generators:
\[
     I=\begin{pmatrix}i&0\\0&-i\end{pmatrix},\qquad
     J=\begin{pmatrix}0&1\\-1&0\end{pmatrix}.
\]
They satisfy $I^2=J^2=-1$ and $JIJ^{-1}=I^{-1}$.
The reduction remains absolutely irreducible: over the
algebraic closure its only possible invariant lines for
$I$ are its two distinct eigenlines, which $J$ exchanges.
Let the original family be $V_0=V$, $V_1=\chi V$.
Both residual representations equal this irreducible
representation because $\omega\equiv1\pmod\pi$.
AC-G38 gives
$T_0=\mathcal O\times_{\mathcal O/\delta}\mathcal O$;
this is local, and it separates the two characteristic-zero
eigensystems.

In the paired family, however, branch $1$ of $V_0$ and
branch $0$ of $V_1$ are the identical representation
$\chi V$.  Every element of $T'$ has equal values at
these two coordinates.  For any nonzero $c\in\mathcal O$,
the element $cP_1$ has values $c$ and $0$ there, and thus
does not belong to $T'$.  In particular $\delta P_1\notin T'$.
No $\tau$ as in (G39.4) exists: $V(\tau)=1$ would give
$(\chi V)(\tau)=\omega\ne1$ whenever $\chi(\tau)=\omega$.
This is a counterexample to the stated Hecke-family
fragment, not a countermodel to the full source data.
\end{proof}

\paragraph{Exact source interface and remaining comparison.}
The source local algebra and generator hypotheses are
\texttt{NData.T\_local} and \texttt{NData.T\_gen} in
\path{Formalism/SharpRestriction/LevelLetterCocycles/DataN.lean}.
The common-residue argument is exactly the proved
\texttt{res\_T\_indep}, followed by \texttt{heckePol\_res}, in
\path{JointConditionsZpowers/ResidualIrreducibleRankR.lean}.
Passage from split unramified Hecke primes to every group
element is the finite-family closedness argument in
\texttt{T\_closed} and
\path{SumRepresentations.charpoly_clause_of_chebotarev}
in \path{ExistsInertiaCharpolyFactorOfTameType.lean}.
The required good residual restriction is the explicit
PointsLemma conclusion already isolated in AC-G37-self.

For coefficient $3$, a prime above $269$ supplies (G39.4)
if the original family is unramified there and the cubic
character has its full inertia image.  The latter holds
for $\chi_{269}$ over the hosts under consideration by
their unramifiedness at $269$ and the tame cubic Kummer
calculation.  The source's precise original-family
unramifiedness clause is \texttt{SumRepresentations.At.G4}
in \path{SumRepresentations.lean}: every original
representation is unramified outside places over the
coefficient prime, \texttt{N.S}, and \texttt{N.Q}.
By \texttt{NData.S\_def}, $S$ consists of the two test
prime sets and the selected auxiliary prime set, with
coefficient-prime places removed.  Thus excluding $269$
from these sets and from the original $Q$ is the exact
condition needed here.  It is not silently assumed for
an arbitrary pre-existing twisted local component.
No such additional condition is needed away from three.

The new result controls the added branch projection even
when the old local congruence algebra $T_0$ is complicated.
To use it in the source inequality, one must still
construct the paired theta lattice $M$, identify each
projected branch $N_i$ integrally with the appropriate
original theta lattice, and identify the selected subspace
and character quotient with the required \texttt{etaJ}
and Selmer coefficient.  Under those identifications,
(G39.10) adds at most one half of the selected rank at
coefficient three.  The normalized global bracket is then
at most $\tfrac12\log3$ if the selected ranks sum to at
most \texttt{dZ0}, by AC-G28-normalization.  This retains
the original index term and the actual rank hypothesis.
The algebra theorem alone neither supplies these
identifications nor reverses the direction of
\texttt{LinkN} to bound a Selmer group.

\paragraph{Dictionary.}
The finite character, its kernel cover, and restriction
after base change use the previously audited correspondence
in SGA~1, Expos\'e~V, \S\S4 and 7.  The fiber-product order
and its projection/intersection indices are proved above;
no new geometric or Hodge-theoretic existence assertion
is attached to that dictionary.
